% Source coverage: physical PDF pp. 485--487 (printed pp. 462--464),
% beginning at the Section 23.7 heading and ending immediately before
% Section 23.8. Numbered equations: (23.7.1)--(23.7.4).
\subsection{Quantum Fluctuations around Extended Field Configurations}
\label{sec:23.7}

Topologically non-trivial four-dimensional field configurations such as
instantons provide only zeroth-order contributions to path integrals.
Now we must consider the effect of quantum fluctuations around these
configurations.

To start in a very general way, consider a set of fields, collectively
called $\phi(x)$, whose dynamics is described by the Euclidean action
$I[\phi]$, and suppose we have a set of configurations $\phi_{\nu,u}$ at
which $I[\phi]$ is stationary, where $\nu$ labels the topological type of
the configuration, and $u$ represents a set of continuous collective
parameters on which the configuration depends. For instance, for
instantons $\nu$ is the winding number, an integer, and $u$ includes the
position and scale as well as its `direction' in the gauge group.
Euclidean path integrals may then be written
\begin{equation}
\int[d\phi]\,\exp(I[\phi])\mathcal O
=\sum_\nu\int du\int_u[d\phi']\,
\exp\bigl(I[\phi_{\nu,u}+\phi']\bigr)\mathcal O,
\tag{23.7.1}\label{eq:23.7.1}
\end{equation}
where the subscript $u$ on the integral over the fluctuation $\phi'$
indicates that we are to integrate only over fluctuations that do not
entail changes in the collective parameters, and $\mathcal O$ stands for
any product of local functions of field operators. Because $I[\phi]$ is
stationary at $\phi=\phi_{\nu,u}$, its expansion to second order in the
fluctuations takes the form
\begin{equation}
I[\phi_{\nu,u}+\phi']\simeq I_\nu
-\frac12\int d^4x\,d^4y\,
K_{xl,ym}(\nu,u)\phi'_l(x)\phi'_m(y),
\tag{23.7.2}\label{eq:23.7.2}
\end{equation}
where $l$ and $m$ here include spin and species indices, and
$I_\nu\equiv I[\phi_{\nu,u}]$ is a function of $\nu$ alone because the action
is supposed to be stationary at all field configurations
$\phi_{\nu,u}$. The integral over the fluctuation $\phi'$ in
Eq.~\eqref{eq:23.7.1} will then yield a sum of terms from contractions
of the fields in $\mathcal O$, times (for real bosonic fields) an
over-all factor $[\operatorname{Det}K(\nu,u)]^{-1/2}$, it being
understood that $K$ now acts only in the subspace of fluctuations that
do not entail changes in the collective parameters $u$. This factor can
be written as a product over the eigenvectors $\lambda_n(\nu,u)$ of the
`matrix' $K(\nu,u)$:
\begin{equation}
\bigl(\operatorname{Det}K(\nu,u)\bigr)^{-1/2}
=\prod_n^{\prime}\lambda_n^{-1/2}(\nu,u),
\tag{23.7.3}\label{eq:23.7.3}
\end{equation}
the prime here indicating that we are to exclude the `zero modes', the
zero eigenvalues of $K$, for which the eigenfunctions correspond to
changes in the collective parameters.

These remarks allow us to refine the results given in
\hyperref[sec:23.5]{Section 23.5} for the coupling-constant dependence of
instanton contributions with different winding numbers. Suppose we
define the fields of our theory in such a way that the action takes the
form $I[\phi,g]=g^{-2}I_1[\phi]$, where $I_1[\phi]$ is independent of
coupling constants. (For instance, as discussed at the end of
\hyperref[sec:15.2]{Section 15.2}, in Yang--Mills theories we would use
as a gauge field the canonically normalized gauge field times a factor
$g$.) Then all eigenvalues of $K$ are proportional to $g^{-2}$, and the
factor \eqref{eq:23.7.3} is proportional to $g$ to a power equal to the
number of non-zero eigenvalues of $K$. This power is of course infinite,
but it can be written as the number of all eigenvalues of $K$, which is
also infinite but is independent of $\nu$, minus the number
$\mathcal N(\nu)$ of zero modes, equal to the number of collective
parameters, which is finite and dependent on $\nu$. We conclude then
that aside from factors that do not depend on $\nu$, the contribution of
fluctuations around configurations of topological type $\nu$ is a factor
with a coupling-constant dependence that is given in terms of the number
$\mathcal N(\nu)$ of collective parameters by
\begin{equation}
g^{-\mathcal N(\nu)}.
\tag{23.7.4}\label{eq:23.7.4}
\end{equation}
This estimate is based on the approximation \eqref{eq:23.7.2}, which
corresponds to a one-loop approximation, so when higher-order terms are
taken into account the factor \eqref{eq:23.7.4} will be multiplied with
a power series in $g$, whose details depend on the operators $\mathcal O$
appearing in the path integral.

Let's see how this applies to instantons. The configuration with $\nu=0$
of course has no collective parameters, and so its contribution to path
integrals is just a power series in $g$. The configuration with $\nu=1$
has four collective parameters giving the spacetime location of the
instanton, one collective parameter giving the scale of the instanton,
and a number $N_1$ of collective parameters corresponding to the
rotations and/or global gauge group transformations that do not leave
the instanton invariant, so (now including the factor $\exp(I_\nu)$ given
by Eq.~\eqref{eq:23.5.19}), the coupling-constant dependence for $\nu=1$
instantons is $g^{-5-N_1}\exp(-8\pi^2/g^2)$. For the $\nu=1$ instanton
\eqref{eq:23.5.12} in an $SU(2)$ Yang--Mills theory, there are three
independent rotations and three independent $SU(2)$ transformations,
but since the instanton is invariant under three combined rotations and
$SU(2)$ transformations, we have $N_1=3$, so fluctuations around the
$\nu=1$ instanton yield a coupling-constant dependence
$g^{-8}\exp(-8\pi^2/g^2)$. In an $SU(3)$ Yang--Mills theory we have three
independent rotations and eight independent $SU(3)$ transformations,
but again the instanton is invariant under three independent combined
rotations and $SU(3)$ transformations in the standard $SU(2)$ subgroup
(say, acting on the first two components of the defining representation
of $SU(3)$), and now also invariant under one additional $SU(3)$
transformation that (like hypercharge) commutes with all generators of
the standard $SU(2)$ subgroup. Hence $N_1=3+8-3-1=7$, and the $\nu=1$
instanton yields a factor proportional to
$g^{-12}\exp(-8\pi^2/g^2)$.

Now suppose that the action $I$ also contains a term
\[
-\int d^4x\int d^4y\,
\bar\psi_l(x)\mathcal K_{lx,my}(\nu,u)\psi_m(y)
\]
involving independent fermion fields $\psi$ and $\bar\psi$. If none of
these fields appear in $\mathcal O$ then the integral over these fields
yields a factor $\operatorname{Det}\mathcal K(\nu,u)$, which vanishes if
$\mathcal K(\nu,u)$ has any zero modes. This result can be simply
understood in terms of the rules for integration over fermionic
parameters. By expanding $\psi(x)$ and $\bar\psi(x)$ in the eigenmodes
of $\mathcal K$, we can write the integral over $\psi(x)$ and
$\bar\psi(x)$ as an integral over the coefficients in these expansions.
The coefficients of the zero modes do not appear in the quadratic
approximation to the action, so for each such fermionic zero mode we have
an integral over a fermionic parameter that does not appear in the
integrand, which vanishes according to the general rules of Section 9.5.
The only terms in an integral over fermionic variables that do not vanish
are those for which the integrand contains a single factor of each
integrated variable. Hence the integral over fermionic fields in
Eq.~\eqref{eq:23.7.1} will not vanish only if there is a single fermionic
field in $\mathcal O$ for each zero mode of $\mathcal K$. For instantons
there are fermionic zero modes whose numbers and chiralities are governed
by index theorems, like the Atiyah--Singer index theorem derived in
\hyperref[sec:22.2]{Section 22.2}, so only certain processes are allowed
for a given winding number. It was on this basis that
't Hooft~\hyperref[ch23-ref-25]{[25]} showed that the baryon- and
lepton-nonconserving effective interaction produced by $\nu=1$
instantons in the electroweak standard model must involve just one of
each lepton flavor.
